\appendix
\section{Certificate formats and independent checking}

\subsection{Sparse polynomials}

A polynomial stores a positive variable count and a sorted list of distinct
exponent vectors with nonzero rational coefficients. Variables are positional.
Exponent vectors have exactly the declared dimension and contain nonnegative
integers. Construction rejects noncanonical encodings instead of silently
merging duplicate entries supplied in a purported certificate.

For example, $x^2-2xy+y^2$ is represented schematically as
\begin{Code}
{
  "nvars": 2,
  "terms": [
    [[0, 2], "1"],
    [[1, 1], "-2"],
    [[2, 0], "1"]
  ]
}
\end{Code}

All fractions in the exported JSON are exact strings such as \code{"1/97"}.
The checker does not parse a decimal into a claim of approximate equality.
An external-facing production parser should additionally bound sizes before
constructing these values.

\subsection{Cone and invariant records}

A cone record contains the target polynomial, ordered lists of nonnegative
and zero hypotheses, and a certificate. Each cone term contains a rational
weight, a square polynomial, and a list of factor indices. An index can occur
more than once. The equality part contains one arbitrary polynomial multiplier
per zero hypothesis. The exact order is part of the certificate's meaning.

An invariant record contains an invariant polynomial, an initial rational state,
and one transition polynomial per state coordinate. The reload checker tests
both base evaluation and exact polynomial substitution. It is independent of
the interpolation and LP discovery functions. It is not independent of the
shared rational-polynomial implementation, a limitation addressed by tests
rather than concealed as formal verification.

The essential cone-checking logic is the following executable pattern:
\begin{Code}
acc = zero_polynomial
for term in certificate.terms:
    reject unless term.weight is rational and nonnegative
    p = term.weight * term.square**2
    for i in term.factors:
        reject unless i is a valid hypothesis index
        p *= nonnegative_hypotheses[i]
    acc += p
for multiplier, equality in paired_ideal_entries:
    acc += multiplier * equality
accept exactly when acc == target
\end{Code}

The actual implementation also checks dimensions, array lengths, canonical
polynomials, and malformed input. Exceptions representing invalid evidence
lead to rejection. Resource-exhaustion behavior is an additional production
requirement, not a theorem proved about the prototype.

\subsection{Horn and box records}

A Horn proof node records its ground conclusion, the rule index or input-fact
marker, a complete ground substitution, and earlier premise-node indices.
The checker verifies the instantiated rule body and head against those exact
nodes. Trimming retains only ancestors of the root, preserving the topological
order. A dependency cycle in the search graph may be legitimate, but the final
proof certificate cannot justify itself through that cycle.

A Bernstein certificate is either a leaf with a box, coordinate degrees, and
all basis coefficients, or a split with an axis, interior rational cut, and
both children. Every basis index must appear exactly once in canonical order.
The checker reconstructs the polynomial at each leaf and checks the complete
box cover recursively.

\section{Reproducing the package}

From the extracted package root, the Python path is:
\begin{Code}
python -m venv .venv
. .venv/bin/activate
python -m pip install -r prototype/requirements.txt
bash scripts/reproduce.sh
\end{Code}

For the exact versions used in this run, substitute
\path{prototype/requirements-tested.txt}. That file is a tested-version
snapshot, not a promise that binary wheels exist for every platform.

Individual actions can be run from \path{prototype/}:
\begin{Code}
python -m pytest -q
python run_benchmarks.py --output ../results
python validate_artifacts.py --results ../results
python export_lean.py --results ../results --output ../lean
\end{Code}

The reproduction script regenerates measurements and certificates; wall-clock
timings naturally vary. It does not download or install Lean. To attempt the
uncompiled replay project on a machine with Lean/Lake and network access:
\begin{Code}
cd lean
lake update
lake exe cache get
lake build
cd ..
python scripts/run_lean_comparisons.py --timeout 30
\end{Code}

The runner preserves compiler output and distinguishes timeout from nonzero
exit. A nonzero compiler exit can be an elaboration or dependency failure,
not necessarily a failed mathematical search. The runtime absence recorded in
the supplied results is intentional and is not counted as either a solve or
a regression.

The article is built with LuaLaTeX. Its main source is
\path{article/forge.tex}, which inputs the numbered chapter files,
\path{horn-table.tex}, and \path{references.tex}. No font files are bundled.
The build helper updates the numerical Horn table from the saved JSON before
running the typesetter.
